\chapter{Theoretical Background}\label{chap:back}
%In this chapter we introduce, in order, the concept of bosonization, the Matsubara formalism, non-hermitian topology and the perturbative \acrlong{rg}.
In this chapter, we introduce the relevant analytical tools and concepts.
Bosonization is discussed in \autoref{chap:bos}, and the Matsubara formalism for the computation of \acrlong{gf}s at finite temperatures in \autoref{chap:gf}.
A brief introduction to \acrlong{nh} topology is given in \autoref{chap:ep}, and two different approaches to the perturbative \acrlong{rg} are introduced in \autoref{chap:rg}.
While occasionally remarks beyond the content of standard references are made, this part primarily aims to provide a reasonably self-contained body of necessary background information.
\section{Bosonization}\label{chap:bos}
At its core, bosonization describes the procedure of re-expressing fermionic operators in terms of bosonic ones \cites[\ssec II]{Bosonization_FiniteT_Bowen_2001}[\ssec 3.I]{Gogolin_1999}, which is essentially an unusual and highly non-trivial change of basis \cites[\spage 34]{Giamarchi_2004}.
The primary goal of this section is to describe the method itself, for which this formal point of view is adequate.
\subsection{Breakdown of Fermi Liquid theory}\label{sec:bos:introduction}
The main motivation for bosonization is the breakdown of Fermi Liquid theory in one dimensional systems \cite{Giamarchi_2004,FermiLiquid_LL_review_Voit_1995}.
In higher dimensions, Fermi Liquid theory very successfully manages to describe interacting electrons as approximately free quasi-particles with a finite lifetime near the Fermi energy \cite{QMB_Mahan_2000,FermiLiquid_review_Neilson_1996,FermiLiquid_LL_review_Voit_1995}.
Since these consist of an electron and density fluctuations, they are also of fermionic nature, and, according to \cite{FermiLiquid_review_Neilson_1996,Giamarchi_2004}, Landau first demonstrated that in $d=3$ dimensions their lifetime scales as $\tau \sim E(k)^{-2}$.
An accessible and detailed derivation of this is given in \cite{FermiLiquid_Vignale_2022}, which also shows that in two dimensions the relation is only slightly modified by logarithmic corrections.
This scaling is important, because it guarantees that as we approach the Fermi level, $\frac{1}{\tau}$ becomes much smaller than the energy $E(k)$, and thus quasi-particles are well defined \cite{Giamarchi_2004,FermiLiquid_LL_review_Voit_1995}.
Consequently, Fermi Liquid theory is a good description of the low energy limit, even for strong interactions \cites[\ssec 1.1]{Giamarchi_2004}[\ssec 6.2]{QMB_Coleman_2015}.
\\
In one dimension, however, this fails and the \trspace{} picture provides an intuitive explanation.
As soon as an interaction is turned on, any particle propagation causes a collective excitation since it is no longer possible to ``go around'' other particles.
But if individual motion is impossible, then the notion of (almost) free quasi-particles can not be an accurate model.
More formally, one can show that in one dimension the susceptibility measuring the linear response to an external potential diverges as $T\rightarrow 0$, suggesting that the ground state is very different from the non-interacting case \cite[\ssec 1.2]{Giamarchi_2004}.
Without giving a detailed account of the derivation, note that a crucial observation is that the so-called \emph{nesting condition} $\epsilon(k+q) = -\epsilon(k)$ is, for $d>1$, generally only satisfied at isolated points.
For $d=1$, the Fermi surface only consists of two points at $k=\pm \kf$, and for $q=2\kf$ the nesting condition is satisfied for an entire interval of momenta within which a linear approximation\footnote{At least assuming inversion symmetry, otherwise the dispersion need not be symmetric. However, note that this is not a requirement for bosonization.} of the dispersion is valid \cite[\spage 8]{Giamarchi_2004}.\\
It should be noted that it is also possible to obtain useful results using perturbation theory \cite[\spage 11]{FermiLiquid_LL_review_Voit_1995}.
In particular, historically the first solution was entirely in terms of fermions \cite{Parquet_Dzyaloshinskii_1974}.
Nevertheless, bosonization provides a rather elegant solution to the problem of interacting electrons in one dimension and leads to the \gls{ll} as a generic low-energy theory in place of the Fermi Liquid.
The concept of the \gls{ll} has proven useful for the quantum impurity problem \cite{QImpurity_review_Furusaki_2005,QImpurity_LL_Rylands_2016}, quantum Hall edge states \cite{QHallEdgeStates_LL_Wen_1990} and, especially in recent years, the helical \gls{ll} has gained a lot of interest \cite{Helical_ChargeSpinDuality_nonequilibrium_Liu_2011,Helical_Edge_PhononBackscat_Budich_2012,Helical_RashbaImp_RG_Crepin_2012,Helical_RashbaNanowires_LowElectronDensity_Schmidt_2016,Helical_RashbaNanowires_Magnetic_Pedder_2016,Helical_Rashba_SpinOrbit_QSH_Stroem_2010,Helical_SyntheticQWire_Japaridze_2013}. % claim for impurity and QHES in BosonReferm p.4
One of the predicted signatures, a re-entrant conductance feature within the helical phase, has even been directly observed in experiment \cite{Helical_Experiment_Nanowire_Heedt_2017}.
Finally, the ``claim to fame'' of bosonization is that often interactions can be taken into account exactly.
This is because some -- but, alas, not all -- products of four fermionic creation and annihilation operators end up being mapped to an expression that is merely quadratic in the new degrees of freedom \cite{Bosonization_Schulz_1995,Giamarchi_2004}.
\subsection{Bosonization Identity}\label{sec:bos:bosonization_identity}
There exists a plethora of introductions to bosonization.
Even within the standard approach \cite{Giamarchi_2004,Gogolin_1999,Fradkin_2013,QMB_Fabrizio_2022,Bosonization_TLL_Schonhammer_1997,Bosonization_GF_LL_Meden_1999,Bosonization_Rao_2001,Bosonization_LowE_1d_fluids_Haldane_1981} the construction varies quite a bit, and famously ``there are as many different notations as there are authors'' \cite[\spage 394]{Giamarchi_2004}.
Some derivations also involve the \acrlong{rg} \cite{Bosonization_Shankar_1994,Bosonization_Schulz_1995} or are even based on Conformal Field Theory \cite{Bosonization_Senechal_2004}.
In this section, we derive the bosonization identity for an arbitrary number of fermionic fields following the comparatively rigorous approach of \cite{Bosonization_DelftSchoeller_1998}.
For our purposes the case of only two fermions would suffice, see e.g. \cite{Giamarchi_2004}.
But the generalization actually \emph{simplifies} the derivation a bit by providing a clearer separation between details that belong to the procedure itself, and details that originate from the model specifics.\\
The starting point is to assume that we are given a set of fermionic operators $c_{k,\ell}$ that satisfy the canonical anti-commutation relations
\begin{align}
\klc[\big]{c_{k,\ell}, c_{k',\ell'}^\dagger} &= \delta_{\ell,\ell'} \delta_{k,k'}\quad\text{and}\quad \klc[\big]{c_{k,\ell}, c_{k',\ell'}} = 0. \label{eqn:bos:ckell_canonical_comm}
\end{align}
At this stage, $\ell$ should be viewed as nothing more than a label distinguishing $M$ different species $\ell=1,\dots,M$.
An essential property of these operators is that, besides the momenta being discrete, we also require them to be unbounded\footnote{This is a subtle but crucial point. For a bounded spectrum one of the key steps of bosonization is not possible, see \autoref{eqn:bos:C_q_op_remnants_def}.}.
For a system size $L$ we assume that the momenta are given by
\begin{align*}
k &= \frac{2\pi}{L} n_k\quad\text{for}\ n_k\in\mathbb{Z}.
\end{align*}
This is rather peculiar, as in the case of a lattice model there would always exist some maximal momentum of the order of $\frac{1}{a}$, or equivalently some maximal $n_k$ of the order of the number of lattice sites $N$ \cite[\spage 8]{Bosonization_Rao_2001}.
For now we tentatively accept that we may add these infinitely many additional momenta to our model without affecting the physics in the low-energy limit \cite{FermiLiquid_LL_review_Voit_1995}, but should anticipate \gls{uv} divergences.\\
%Note that thus far we have not assumed anything about the dispersion relation of the system $\epsilon(k)$, although we implicitly assume the existence of a ground state.
For convenience of notation let $\intx{ }{ }{ }{x} \equiv \intx{-L/2}{L/2}{ }{x}$ and $\sum_k \equiv \sum_{n_k\in\mathbb{Z}}$.
We define field operators
\begin{align}
\psi_\ell(x) &= \frac{1}{\sqrt{L}} \sum_k \powe{\i kx} c_{k,\ell} \label{eqn:bos:psi_ell_from_c_k_ell}
\end{align}
where $\intx{ }{ }{\powe{\i kx}}{x}=L\delta_{k,0}$ implies the inverse relation
\begin{align}
c_{k,\ell} &= \frac{1}{\sqrt{L}} \intx{ }{ }{\powe{-\i kx} \psi_\ell(x)}{x}. \label{eqn:bos:c_k_ell_from_psi_ell}
\end{align}
Using \autoref{eqn:bos:ckell_canonical_comm} and \autoref{eqn:bos:psi_ell_from_c_k_ell}, the anti-commutation relations for the field operators become\footnote{Since $f(x) = \sum_{n\in\mathbb{Z}} \powe{\i nx}$ is $2\pi$-periodic it can be interpreted as a Fourier series with coefficients $c_n=1$. Due to $c_n = \frac{1}{2\pi} \intx{p}{p+2\pi}{\powe{-\i nx} f(x)}{x}$ for any $p\in\mathbb{R}$ we must have $f(x) = 2\pi \sum_{m\in\mathbb{Z}} \delta(x - 2\pi m)$, see \cite[\spages 30--32]{GelfandShilov_1964} for details.} % Deitmar and Echterhoff "Principles of Harmonic Analysis" Thm. 3.6.3. OR I. M. Gel’fand, G. E. Shilov, Generalized Functions, Vol. 1, Academic Press, New York 1964 p. 331 (OR p.31, (2.4.2))
\begin{align*}
\klc*{\psi_\ell(x), \psi_{\ell'}^\dagger(x')} &= \delta_{\ell,\ell'} \frac{1}{L} \sum_{n_k\in\mathbb{Z}} \powe{\i \frac{2\pi}{L} n_k (x-x')} = \delta_{\ell,\ell'} \sum_{m\in\mathbb{Z}} \delta\kla*{x-x'-mL}.
\end{align*}% = \frac{1}{L} \sum_{k,k'} \powe{\i kx} \powe{-\i k'x'} \klc*{c_{k,\ell},c_{k',\ell'}^\dagger}
This serves as a periodic generalization of -- and simplifies to -- the usual result in the case of $x,x'\in \klb*{-\frac{L}{2},\frac{L}{2}}$ \cite[\spage 10]{Bosonization_DelftSchoeller_1998}.
\myparagraph{Vacuum state}
Before we begin with the proper construction of bosonization, we have to think about the notion of a ground state in our system.
A rather natural definition is given by the properties
\begin{align*}
c_{k>0,\ell} \Ket{\textbf{0}}_0 &=0,\quad\text{and}\quad c_{k\le 0,\ell}^\dagger \Ket{\textbf{0}}_0 = 0.
\end{align*}
One can think of this as the Fermi sea where every state up to and including $k=0$ is occupied for all labels $\ell$, and all higher momentum states are empty \cite{Bosonization_DelftSchoeller_1998}.
Although we refer to this as the ground state, note that it is actually not necessarily the state of lowest energy.
However, it is the reference point with respect to which we measure everything else.
In principle this choice is somewhat arbitrary due to the unbounded spectrum; we could just as well choose the state $c_{0,1} \Ket{\textbf{0}}_0$ containing one fewer fermion with label $\ell=1$ and momentum $k=0$.
More generally, there exist infinitely many $\textbf{N}$-particle ground states that we define by adding (or removing) the respective number of fermions of each label contained in the vector $\textbf{N}=(N_1,\dots,N_M)$ to (or from) $\Ket{\textbf{0}}_0$.
To disambiguate the phase, we specify an order in the same way as \cite{Bosonization_DelftSchoeller_1998}.
Identifying $c_{n_k,\ell} \equiv c_{k,\ell}$ for convenience, we define
\begin{align*}
\Ket{\textbf{N}}_0 &= (C_1)^{N_1} \dots (C_M)^{N_M} \Ket{\textbf{0}}_0\quad\text{with} \quad (C_\ell)^{N_\ell} = \begin{cases} c_{N_\ell,\ell}^\dagger \dots c_{1,\ell}^\dagger & \text{for}\ N_\ell>0,\\
1 & \text{for}\ N_\ell = 0,\\ c_{N_\ell+1,\ell}\dots c_{0,\ell} & \text{for}\ N_\ell < 0. \end{cases}
\end{align*}
As an illustrative example, this translates to $\Ket{(0,3,-2)}_0 = c_{3,2}^\dagger c_{2,2}^\dagger c_{1,2}^\dagger c_{-1,3} c_{0,3} \Ket{(0,0,0)}_0$.\\
%The essence of bosonization can then be expressed as the following claim: any state created by the fermionic $c_{k,\ell}$ acting on the Fermi sea $\Ket{\textbf{0}}_0$ can also be obtained by a combination of particle-hole excitations and raising or lowering operators that change the total number of fermions.\\
%More formally, the goal is to show that the space of states can be reorganized as a direct sum $\bigoplus_\textbf{N} \mathcal{H}_\textbf{N}$ where $\mathcal{H}_\textbf{N}$ denotes the subspace containing a fixed number of particles of each label corresponding to the vector $\textbf{N}$.\\
%Within the Hilbert space $\mathcal{H}_\textbf{N}$ of a fixed number of particles $\textbf{N}$ for the respective labels all excitations must preserve the total number of particles and must therefore be of bosonic nature, i.e. come in the form of particle-hole excitations.
Let us now consider the density operator in \trspace{},
$\rho_\ell(x) = \varnormord{\psi_\ell^\dagger(x) \psi_\ell(x)}$, where the notation $\varnormord{A} \equiv A - \bra*{A}_0$ serves as a regularization with respect to the Fermi sea $\Ket{\textbf{0}}_0$.
The Fourier components of this operator take the form of particle hole excitations, because
\begin{align*}
\rho_{q,\ell} &= \intx{ }{ }{\powe{-\i qx}\rho_\ell(x)}{x} = \frac{1}{L} \sum_{k,k'} \intx{ }{ }{ \powe{-\i qx} \powe{-\i kx} \powe{\i k'x} \varnormord{c_{k,\ell}^\dagger c_{k',\ell}}}{x} = \sum_k \varnormord{c_{k,\ell}^\dagger c_{k+q,\ell}}.
\end{align*}
Note that the regularization is only necessary for $q=0$, since otherwise $\bra[\big]{c_{k+q,\ell}^\dagger c_{k,\ell}}_0 = 0$ by construction. 
This case corresponds to the number operator
\begin{align*}
N_\ell &= \rho_{q=0,\ell} = \sum_k \varnormord{c_{k,\ell}^\dagger c_{k,\ell}}
\end{align*}
counting the number of fermions with label $\ell$ relative to the Fermi sea.
Somewhat unusually, here the eigenvalues $\mathcal{N}_\ell$ can also be negative, which is a direct consequence of $c_{k<0,\ell} \Ket{\textbf{0}}_0 \neq 0$ \cite{Bosonization_DelftSchoeller_1998}.
The number operator admits the usual commutation relations
\begin{align*}
\klb*{N_\ell, c_{k,\ell'}} &= \delta_{\ell,\ell'} \klb*{c_{k,\ell}^\dagger c_{k,\ell}, c_{k,\ell}} = -\delta_{\ell,\ell'} c_{k,\ell}\quad \text{and} \quad \klb*{N_\ell, c_{k,\ell'}^\dagger} = \delta_{\ell,\ell'} c_{k,\ell}^\dagger.
\end{align*}
\vspace*{-1em}
\myparagraph{Bosonic operators}
To analyse the commutation relations of the $\rho_{q,\ell}$ for $q\neq 0$, note that due to $\rho_{q,\ell}^\dagger = \rho_{-q,\ell}$ we only need to consider
\begin{align*}
\rho_{q,\ell} \rho_{q',\ell'} &= \sum_{k,k'}c_{k,\ell}^\dagger c_{k+q,\ell} c_{k',\ell'}^\dagger c_{k'+q',\ell'} = \sum_{k,k'}c_{k,\ell}^\dagger \klc*{\delta_{\ell,\ell'} \delta_{k+q,k'} - c_{k',\ell'}^\dagger c_{k+q,\ell} } c_{k'+q',\ell'}\\
&= -\sum_{k,k'} c_{k',\ell'}^\dagger c_{k,\ell}^\dagger c_{k'+q',\ell'} c_{k+q,\ell} + \delta_{\ell,\ell'} \sum_k c_{k,\ell}^\dagger c_{k+q+q',\ell}\\
&= \rho_{q',\ell'}\rho_{q,\ell} + \delta_{\ell,\ell'} \sum_k c_{k,\ell}^\dagger c_{k+q+q',\ell} - \delta_{\ell,\ell'} \sum_k c_{k-q',\ell}^\dagger c_{k+q,\ell}.
\end{align*}
At this point, one can split the operators in the summation into their regularized version and the expectation value.
This only makes a difference for $q'=-q$, so we restrict ourselves to that case.
Now consider for a moment a bounded spectrum, i.e. $n_k=n_\text{min}, \dots, 0, \dots, n_\text{max}$.
Then the commutator becomes \cites[\ssec 2.1]{Giamarchi_2004}{Bosonization_DelftSchoeller_1998}
\begin{align*}
\klb*{\rho_{q,\ell},\rho_{-q,\ell}} &= \sum_{n_{k}=n_\text{min}}^{n_\text{max}} \klb*{\varnormord{c_{k,\ell}^\dagger c_{k,\ell}} + \bra*{c_{k,\ell}^\dagger c_{k,\ell}}_0 - \varnormord{c_{k+q,\ell}^\dagger c_{k+q,\ell}} - \bra*{c_{k+q,\ell}^\dagger c_{k+q,\ell}}_0}\\
&= C_{q} + \sum_{n_{k}=n_\text{min}}^{0} - \sum_{n_{k}=n_\text{min}}^{-n_q} = C_{q} + n_q
\end{align*}
where the remaining operator components are absorbed in
\begin{align*}
C_{q} &= \sum_{n_{k}=n_\text{min}}^{n_\text{max}} \varnormord{c_{k,\ell}^\dagger c_{k,\ell}} - \sum_{n_{k}=n_\text{min}}^{n_\text{max}} \varnormord{c_{k+q,\ell}^\dagger c_{k+q,\ell}}. \numberthis \label{eqn:bos:C_q_op_remnants_def}
\end{align*}
Since these operators are regularized, we may perform an index shift $k\rightarrow k-q$ in the second summation \cite[\spages 32--33]{Giamarchi_2004}.
However, unless $n_\text{min}=n_\text{max}\rightarrow\infty$, there is only partial cancellation and the commutation relations of the $\rho_{q,\ell}$ remain operator valued and thus too complicated to proceed.
For an unbounded spectrum, the index shift instead makes the two series identical, giving $C_q = 0$ and therefore
\begin{align*}
\klb*{\rho_{q,\ell}, \rho_{q',\ell'}} &= \delta_{\ell,\ell'} \delta_{q,-q'} n_q. %\label{eqn:bos:rho_qell_commutator}
\end{align*}
Now we are finally in a position to define, for $q>0$, the set of bosonic operators
\begin{align}
b_{q,\ell} &= \frac{-\i}{\sqrt{n_q}} \rho_{q,\ell} = \frac{-\i}{\sqrt{n_q}} \sum_k c_{k-q,\ell}^\dagger c_{k,\ell} \label{eqn:bos:bq_ell_def}
\end{align}
and $b_{q,\ell}^\dagger \equiv \kla[\big]{b_{q,\ell}}^\dagger$ \cite{Bosonization_DelftSchoeller_1998}.
These operators then fulfil the canonical commutation relations\footnote{The second relation only holds, because by definition $q,q'>0$, and thus $\delta_{q,-q'}$ vanishes.}
\begin{align*}
\klb[\big]{b_{q,\ell}, b_{q',\ell'}^\dagger} &= \frac{1}{\sqrt{n_q n_{q'}}} \klb[\big]{\rho_{q,\ell}, \rho_{-q',\ell'}} = \delta_{\ell,\ell'} \delta_{q,q'} \quad \text{and} \quad \klb[\big]{b_{q,\ell}, b_{q',\ell'}} = 0.
\end{align*}
The $\textbf{N}$-particle ground state defined above is compatible with the notion of a ground state implied by the particle-hole excitations, since $b_{q,\ell} \Ket{\textbf{N}}_0 = 0$.
To understand why this is the case, we need to consider how $b_{q,\ell}$ acts on an arbitrary state. 
The label $\ell$ does not play a role here, so we drop it and also only focus on a single species for notational convenience.
For $q>0$, the operator $c_{k-q}^\dagger c_{k}$ attempts to shift a fermion with momentum $k$ to the lower momentum $k-q$.
If $k$ is not occupied or $k-q$ already is occupied, then the state is destroyed.
Since in $\Ket{N}_0$ all momenta up to $n_N$ are filled without gaps, the contributions to $b_q\Ket{N}_0$ can only come from $n_k \le n_N$.
But we can not move these fermions to lower states, because they are already occupied, formally $c_{k-q}^\dagger \Ket{N}_0 = 0$.
Thus, $c_{k-q}^\dagger c_{k}$ destroys $\Ket{N}_0$ for any $k$, and it follows that $b_q\Ket{N}_0=0$.
\myparagraph{Operator identity}
Now that we have a proper ground state for every subspace, we can define the excited states in a very natural way via
\begin{align}
\Ket{\textbf{N},\klc{\textbf{m}_q}} &= \prod_\ell \prod_{q>0} \frac{1}{\sqrt{m_{\ell,q}}} \kla[\big]{b_q^\dagger}^{m_{\ell,q}} \Ket{\textbf{N}}_0
\end{align}
where $\klc{\textbf{m}}_q$ is a sequence of vectors $\textbf{m}_q = (m_{1,q}, \dots, m_{M,q}) \in \mathbb{N}_0^M$ for every $q>0$.
There is no need to specify a certain order here, because the bosonic operators commute.
The set of all these states forms the basis of the Hilbert space $\mathcal{H}_\textbf{N}$ \cite{Bosonization_DelftSchoeller_1998}.\\
The last ingredient we need is a way to switch from one subspace to another.
This is made possible by defining Klein factors $\eta_\ell$, which remove a fermion with label $\ell$, and the respective counterpart $\eta_\ell^\dagger$.
They are assumed to commute with the bosonic operators \cite[\ssec 4.F]{Bosonization_DelftSchoeller_1998},
\begin{align*}
\klb[\big]{\eta_\ell, b_{q,\ell'}} &= 0 = \klb[\big]{\eta_\ell, b_{q,\ell'}^\dagger},
\end{align*}
and the removal or addition of fermions is defined via
\begin{align*}
\eta_\ell \Ket{\textbf{N}}_0 &= T_\ell \Ket{(N_1, \dots, N_\ell - 1, \dots, N_M)}_0,\\
\eta_\ell^\dagger \Ket{\textbf{N}}_0 &= T_\ell \Ket{(N_1, \dots, N_\ell + 1, \dots, N_M)}_0.
\end{align*}
Here $T_\ell = (-1)^{N_1} \dots (-1)^{N_{\ell-1}}$ is a phase factor that reflects the fermionic nature of the Klein factors.
Consequently the order matters, and for $\ell\neq \ell'$ one finds $\eta_\ell \eta_{\ell'} \Ket{\textbf{N}}_0 = - \eta_{\ell'} \eta_\ell \Ket{\textbf{N}}_0$.
More explicit representations are also possible, see  e.g. \cites[\sapp B]{Bosonization_TLL_Schonhammer_1997}{Bosonization_DelftSchoeller_1998}, but for our purposes the above is sufficient.\\
Since the Klein factors commute with the bosons, we only need to consider the action on the $\textbf{N}$-particle ground states.
It is then a slightly tedious but elementary task to show that $\eta_\ell\eta_\ell^\dagger = \eta_\ell^\dagger\eta_\ell = 1$, as well as\footnote{The same can be found in \cite{Bosonization_DelftSchoeller_1998}, but without derivation, so we want to at least sketch the procedure.
As a simple illustrative example note that for $\ell < \ell'$ applying $\eta_\ell \eta_{\ell'}$ gives $T_\ell T_{\ell'} = +1$ while for $\ell > \ell'$ we instead find $T_\ell T_{\ell'} = -1$; therefore $\klc*{\eta_\ell, \eta_{\ell'}} = 0$ for $\ell\neq \ell'$. 
The other relations follow in a similar fashion.}
\begin{align*}
\klc[\big]{\eta_\ell, \eta_{\ell'}^\dagger} &= 2\delta_{\ell,\ell'}, \quad \klc[\big]{\eta_\ell, \eta_{\ell'}} = 2\delta_{\ell,\ell'} \eta_\ell^2 \quad \text{and} \quad \klb[\big]{N_\ell, \eta_{\ell'}} = -\delta_{\ell,\ell'} \eta_\ell .
\end{align*}
With this, we are ready to define the bosonic field operators and derive the bosonization identity.
Let $\varphi_\ell\eplus \equiv \kla[\big]{\varphi_\ell\eminus}^\dagger$ where
\begin{align}
\varphi_\ell\eminus(x) &= -\sum_{q>0} \frac{1}{\sqrt{n_q}} \powe{-\alpha \frac{q}{2}} \powe{\i qx} b_{q,\ell} \label{eqn:bos:varphi_ell_minus_def}
\end{align}
and also define the hermitian combination $\varphi_\ell\equiv \varphi_\ell\eplus + \varphi_\ell\eminus$.
Here, $\alpha$ is a regularization necessary to avoid \gls{uv} divergences and can be thought of as generating a finite bandwidth $\Lambda= \frac{1}{\alpha}$.
From a physical point of view, the introduction of a short-distance cut-off to artificially limit the bandwidth in this way should have a negligible impact on the low-energy limit.
Recall that we expect to encounter an issue like this due to the artificial extension of the spectrum; for a lattice model the bandwidth is $\Lambda = \frac{\pi}{a}$, but generally bosonization is only valid in a smaller region of \tpspace{} due to the linearization.
Thus, usually $\alpha>a$, and in such a realistic setting it is then neither necessary nor useful to try and set $\alpha\rightarrow 0$ at some point.
There is a caveat to this line of reasoning though, because in the correlation functions, $\alpha$ really has to be a small quantity in order to satisfy certain fundamental requirements, see the discussion in \autoref{sec:ell:gf:preliminaries}.\\
Going back to the Fourier components of the density operator, we now find
\begin{align*}
\rho_\ell(x) &= \frac{1}{L} \sum_{q\neq 0} \powe{\i qx} \rho_{q,\ell} = \frac{1}{L} \sum_{q> 0} \powe{\i qx} \i\sqrt{n_q} b_{q,\ell} + \hc + \frac{1}{L} \rho_{q=0,\ell}\\
&= \frac{\partial_x}{2\pi} \sum_{q>0} \frac{1}{\sqrt{n_q}} \powe{\i qx} b_{q,\ell} + \hc + \frac{1}{L} N_\ell = - \frac{1}{2\pi} \partial_x \varphi_\ell(x) + \frac{1}{L} N_\ell \numberthis \label{eqn:bos:density operator boson from definition}
\end{align*}
where the last equality is only correct in the limit $\alpha\rightarrow 0$.
Using the definition of the $b_{q,\ell}$ in \autoref{eqn:bos:bq_ell_def}, we can compute\footnote{$\klb*{AB,C}=A\klc*{B,C} - \klc*{A,C}B$}
\begin{align*}
\klb*{b_{q,\ell}, \psi_{\ell'}(x)} &= \frac{-\i}{\sqrt{n_q}} \frac{1}{\sqrt{L}} \sum_{k,k'} \powe{\i k'x}\klb*{c_{k-q,\ell}^\dagger c_{k,\ell}, c_{k',\ell'}} = \delta_{\ell,\ell'} \frac{\i}{\sqrt{n_q}} \frac{1}{\sqrt{L}} \sum_{k} \powe{\i (k-q) x} c_{k,\ell}\\
%&= \delta_{\ell,\ell'}\frac{-\i}{\sqrt{n_q}} \powe{\i qx} \psi_\ell(x) \\
&= \delta_{\ell,\ell'} a_{q,\ell}(x)\psi_\ell(x)
\end{align*}
where $a_{q,\ell}(x) = \frac{\i}{\sqrt{n_q}} \powe{-\i qx}$.
Since $b_{q,\ell} \Ket{\textbf{N}}_0 = 0$, we thus find that $\psi_\ell(x) \Ket{\textbf{N}}_0$ must be a coherent state because
\begin{align*}
b_{q,\ell} \psi_\ell(x) \Ket{\textbf{N}}_0 &= \klb*{b_{q,\ell}, \psi_\ell(x)} \Ket{\textbf{N}}_0 = a_{q,\ell}(x) \psi_\ell(x) \Ket{\textbf{N}}_0,
\end{align*}
and can therefore be written as\footnote{$\klb*{A, \powe{B}} = C \powe{B}$ for $C=\klb*{A,B}\in\mathbb{C}$ implies $\psi_\ell \sim \powe{B}$ with $\klb*{b_{q,\ell}, B} = a_{q,\ell}(x)$}
\begin{align*}
\psi_\ell(x) \Ket{\textbf{N}}_0 &= w_\ell(x) \powe{\sum_{q> 0} a_{q,\ell}(x) b_{q,\ell}^\dagger} \eta_\ell \Ket{\textbf{N}}_0 = w_\ell(x) \powe{-\i \varphi_\ell\eplus(x)} \eta_\ell \Ket{\textbf{N}}_0.
\end{align*}
%This is already a remarkable result, but let us delay its interpretation until having obtained the final identity.
Due to $b_q\Ket{N_+,N_-}_0=0$, we can rewrite this in a way that is more symmetric -- and with the power of hindsight more useful -- as a preliminary operator identity
\begin{align*}
\tilde{\psi}_\ell(x) &= \eta_\ell w_\ell(x) \normord{\powe{-\i\varphi_\ell(x)}}
\end{align*}
for which $\tilde{\psi}_\ell \equiv \psi_\ell$ when acting on any ground state $\Ket{N_+,N_-}_0$ \cite[\spage 19]{Bosonization_DelftSchoeller_1998}.
The factor $w_\ell(x)$ still has to be determined, but let us first discuss another subtlety:
Above, the notation $\normord{(\dots)}$ refers to normal-ordering with respect to the bosonic operators in the sense of putting all $b_{q,\ell}^\dagger$ to the left of an expression, e.g. $\normord{b_{q,\ell} b_{q,\ell}^\dagger b_{q,\ell} b_{q,\ell}^\dagger} = \kla{b_{q,\ell}^\dagger}^2 b_{q,\ell}^2$.
Note that despite the claim to the contrary in \cite[\spage 11]{Bosonization_DelftSchoeller_1998}, this is generally \emph{not} equivalent to the regularization $\varnormord{(\dots)}$ that is also sometimes referred to as normal-ordering.
For this specific example one would instead get $\varnormord{b_{q,\ell} b_{q,\ell}^\dagger b_{q,\ell} b_{q,\ell}^\dagger} = \kla{b_{q,\ell}^\dagger}^2 b_{q,\ell}^2 + 3b_{q,\ell}^\dagger b_{q,\ell}$.
The only non-trivial case where there is no difference is that of a bilinear expression, see also \cite[\spage 95]{Cardy_1996}.\\
We can determine the normalization $w_\ell(x)$ by evaluating ${}_0\brc[\big]{\textbf{N}}{\eta_\ell^\dagger \psi_\ell(x)}{\textbf{N}}_0$ in two different ways \cite[\spage 20]{Bosonization_DelftSchoeller_1998}.
First, from the definition we have
\begin{align*}
\frac{1}{\sqrt{L}} \sum_k \powe{\i kx} {}_0\brc*{(N_1,\dots,N_\ell-1,\dots,N_M}{T_\ell c_{k,\ell}}{\textbf{N}}_0 &= \frac{1}{\sqrt{L}} \powe{\i k^\ast x}
\end{align*}
where $k^\ast$ is the momentum of the highest occupied state, i.e. $k^\ast = \frac{2\pi}{L} N_\ell$.
For every other $k$ the annihilation operator either destroys the state, or produces one that is orthogonal to ${}_0\Bra{\textbf{N}}\eta_\ell$.
Second, inserting our preliminary operator identity gives
\begin{align*}
{}_0\brc[\big]{\textbf{N}}{ \eta_\ell^\dagger \eta_\ell w_\ell(x) \powe{-\i\varphi_\ell\eplus(x)} \powe{-\i\varphi_\ell\eminus(x)}}{\textbf{N}}_0 &= {}_0\brc[\big]{\textbf{N}}{ w_\ell(x) \powe{-\i\varphi_\ell\eplus(x)} }{\textbf{N}}_0.
\end{align*}
Assuming that $\Ket{\textbf{N}}_0$ is an eigenstate of $w_\ell$, we can then pull out its eigenvalue and act with the remaining exponential operator to the left.
This leaves just the eigenvalue of $w_\ell$ and our assumption is therefore immediately verified as we find $w_\ell(x) = \frac{1}{\sqrt{L}} \powe{\i \frac{2\pi}{L} N_\ell x}$.\\
It remains to be shown that the operator identity also holds when applied to an arbitrary state $\Ket{\textbf{N},\klc*{\textbf{m}_q}}$.
Before we prove that this is indeed the case, note that
\begin{align*}
\psi_\ell(x) b_{q,\ell'}^\dagger &= b_{q,\ell'}^\dagger \psi_\ell(x) + \klb*{\psi_\ell(x), b_{q,\ell'}^\dagger} = \kla*{b_{q,\ell'}^\dagger - \delta_{\ell,\ell'} a_{q,\ell}^\ast(x)} \psi_\ell(x)
\end{align*}
which can be readily generalized to hold for arbitrary powers of $b_q^\dagger$ and consequently any function that admits a series representation.
Here, we merely need that
\begin{align*}
\psi_\ell(x) \Ket{\textbf{N},\klc{\textbf{m}_q}} &= \psi_\ell(x) \prod_{\ell'} \prod_{q> 0} \frac{1}{\sqrt{m_{q,\ell'}!}} \kla*{b_{q,\ell'}^\dagger}^{m_{q,\ell'}} \Ket{\textbf{N}}_0 \\
&= \klb[\bigg]{\prod_{\ell'} \prod_{q> 0} \frac{1}{\sqrt{m_{q,\ell'}!}} \kla*{b_{q,\ell'}^\dagger - \delta_{\ell,\ell'} a_{q,\ell}^\ast(x)}^{m_{q,\ell'}}} \psi_\ell(x) \Ket{\textbf{N}}_0.
\end{align*}
Since we already know that $\psi_\ell(x) \Ket{\textbf{N}}_0 =\tilde{\psi}_\ell(x) \Ket{\textbf{N}}_0$, we can rewrite this as
\begin{align*}
\tilde{\psi}_\ell(x) \powe{\i\varphi_\ell\eminus(x)} \klb[\bigg]{\prod_{\ell'} \prod_{q> 0} \frac{1}{\sqrt{m_{q,\ell'}!}} \kla*{b_{q,\ell'}^\dagger - \delta_{\ell,\ell'} a_{q,\ell}^\ast(x)}^{m_{q,\ell'}}} \powe{-\i\varphi_\ell\eminus(x)} \Ket{\textbf{N}}_0.
\end{align*}
Repeatedly inserting the identity\footnote{$\powe{A}B\powe{-A} = B + \klb*{A,B} + \frac{1}{2} \klb*{A, \klb*{A,B}} + \dots$, see \cite[\spage 48]{Bosonization_DelftSchoeller_1998}}
%\begin{align}
%\powe{\i\varphi_\ell\eminus(x)} \kla*{b_{q,\ell'}^\dagger - \delta_{\ell,\ell'} a_{q,\ell}^\ast(x)} \powe{-\i\varphi_\ell\eminus(x)} &= \kla*{b_{q,\ell'}^\dagger - \delta_{\ell,\ell'} a_{q,\ell}^\ast(x)} + \klb*{\i\varphi_\ell\eminus(x), b_{q,\ell'}^\dagger} = b_{q,\ell'}^\dagger
%\end{align}
\begin{align*}
\powe{\i\varphi_\ell\eminus(x)} \kla*{b_{q,\ell}^\dagger - a_{q,\ell}^\ast(x)} \powe{-\i\varphi_\ell\eminus(x)} &= \kla*{b_{q,\ell}^\dagger - a_{q,\ell}^\ast(x)} + \klb*{\i\varphi_\ell\eminus(x), b_{q,\ell}^\dagger} = b_{q,\ell}^\dagger
\end{align*}
then allows us to simplify the result to
\begin{align*}
\psi_\ell(x) \Ket{\textbf{N},\klc*{\textbf{m}_q}} &= \tilde{\psi}_\ell(x) \prod_{\ell'} \prod_{q> 0} \frac{1}{\sqrt{m_{q,\ell'}!}} \kla*{b_{q,\ell'}^\dagger}^{m_{q,\ell'}} \Ket{\textbf{N}}_0 = \tilde{\psi}_\ell(x) \Ket{\textbf{N}, \klc*{\textbf{m}_q}}.
\end{align*}
This concludes the proof that the operator identity $\tilde{\psi}_\ell \equiv \psi_\ell$ is indeed valid on the entire space of states, and the relation
\begin{align}
\psi_\ell(x) &= \frac{\eta_\ell}{\sqrt{L}} \powe{\i \frac{2\pi}{L} N_\ell x} \normord{\powe{-\i\varphi_\ell(x)}} \label{eqn:bos:bosonization identity}
\end{align}
is referred to as the \emph{bosonization identity} \cite{Giamarchi_2004,FermiLiquid_LL_review_Voit_1995,Bosonization_DelftSchoeller_1998}.
Exponentials of bosonic fields such as $\powe{m\varphi_\ell(x)}$ are known as \emph{vertex operators} \cite[\ssec 9]{Bosonization_DelftSchoeller_1998}.\\
The derivation is actually not quite complete, because we have only shown the identity to be valid on all the states $\Ket{\textbf{N}, \klc*{\textbf{m}_q}}$.
What we have not proven, is that the collection of these states is sufficient to cover the Hilbert space spanned by the fermionic operators.
The key insight is that this combinatoric problem must be independent of the specific form of the Hamiltonian.
Therefore, it suffices to show that the partition function calculated in the fermionic basis matches the bosonic one for any -- conveniently chosen -- sample system.
An accessible proof can be found in \cites[\sapp B]{Giamarchi_2004}[\sapp B]{Bosonization_DelftSchoeller_1998}, but it is still somewhat lengthy.%\footnote{And the Giamarchi concedes that it does not add much to the physics \cite[p.35]{Giamarchi_2004}.}.
\subsection{Application to common terms}\label{sec:bos:common_terms}
In this section, we apply the bosonization identity, \autoref{eqn:bos:bosonization identity}, to common types of terms.
Note that such terms also appear in the treatment of our model in \autoref{chap:ell}, although with some additional complications.
One of the main tricks used repeatedly is the special case of the Baker-Campbell-Hausdorff relation
\begin{align}
\powe{A} \powe{B} &= \powe{A+B} \powe{\frac{1}{2}\klb*{A,B}} \label{eqn:bos:baker_campbell}
\end{align}
valid if $C=\klb*{A,B}$ commutes with $A$ and $B$ \cite[\spage 48]{Bosonization_DelftSchoeller_1998}.
Applying the formula twice also yields the useful identity $\powe{A} \powe{B} = \powe{B}\powe{A} \powe{\klb*{A,B}}$.
In our case, $C$ is always a simple $\mathbb{C}$-number, so this requirement is satisfied.
More explicitly, using $\sum_{n\ge 1} \frac{1}{n} \powe{-nx} = - \log\kla*{1 - \powe{-x}}$ we find
\begin{align*}
\klb*{\varphi_\ell\eplus(x), \varphi_{\ell'}\eminus(x')} &=  \sum_{q,q'>0} \frac{1}{\sqrt{n_qn_{q'}}} \powe{-\alpha \frac{q+q'}{2}} \powe{-\i qx} \powe{\i q'x'} \klb*{b_{q,\ell}^\dagger, b_{q',\ell'}}\\
&= -\delta_{\ell,\ell'} \sum_{q>0} \frac{1}{n_q} \powe{-\alpha q} \powe{-\i q(x-x')} = \delta_{\ell,\ell'} \log\klb*{1 - \powe{-\frac{2\pi}{L} \kla*{\alpha + \i (x-x')}}}. \numberthis \label{eqn:bos:commutator_varphi_varphi_exact}
\end{align*}
In the thermodynamic limit, this commutator asymptotically behaves as
\begin{align}
\klb*{\varphi_\ell\eplus(x), \varphi_{\ell'}\eminus(x')} &\overset{L\rightarrow\infty}{\sim} \delta_{\ell,\ell'} \log\klb*{\frac{2\pi}{L} \kla*{\alpha + \i (x-x')}}. \label{eqn:bos:commutator_varphi_varphi_tdl}
\end{align}
Let us start with the density operator, but include an offset $\epsilon$ in one of the field operators.
This so-called \emph{point-splitting} \cite[\sapp G]{Bosonization_DelftSchoeller_1998} avoids divergences and yields
\begin{align*}
\psi_\ell^\dagger(x) \psi_\ell(x+\epsilon) &= \normord{\powe{\i\varphi_\ell(x)}} \powe{-\i\frac{2\pi}{L} N_\ell x} \frac{\eta_\ell^\dagger \eta_\ell}{L} \powe{\i\frac{2\pi}{L} N_\ell (x+\epsilon)} \normord{\powe{-\i\varphi_\ell(x+\epsilon)}} \\
&= \frac{1}{L} \powe{\i\frac{2\pi}{L} N_\ell \epsilon} \powe{\i\varphi_\ell\eplus(x)} \powe{\i\varphi_\ell\eminus(x)} \powe{-\i\varphi_\ell\eplus(x+\epsilon)} \powe{-\i\varphi_\ell\eminus(x+\epsilon)} \\
&= \frac{1}{L} \powe{\i\frac{2\pi}{L} N_\ell \epsilon} \normord{\powe{\i\varphi_\ell(x)} \powe{-\i\varphi_\ell(x+\epsilon)}} \powe{\klb*{\i \varphi_\ell\eminus(x), -\i\varphi_\ell\eplus(x+\epsilon)}} \\
%&= \frac{1}{L} \powe{\i\frac{2\pi}{L} N_\ell \epsilon} \normord{\powe{-\i \sum_{n\ge 1} \epsilon^n \frac{1}{n!} \partial_x^n \varphi_\ell(x)}} \klb*{1 - \powe{-\frac{2\pi}{L} \i\epsilon}}^{-1}\\
&= \frac{1}{2\pi\i\epsilon} \frac{1}{1 - \frac{\pi\i\epsilon}{L} + \dots} \normord{\klb*{1 - \i\epsilon\partial_x\varphi_\ell(x) + \mathcal{O}\kla*{\epsilon^2}}} \klb[\Big]{1 + \i\frac{2\pi}{L} N_\ell \epsilon}. \numberthis \label{eqn:bos:psi dagger psi epsilon}
\end{align*}
The regularization included in the definition of the density operator takes care of the divergent constant, and in the limit we find
\begin{align*}
\rho_\ell(x) &= \lim_{\epsilon\rightarrow 0} \varnormord{\psi_\ell^\dagger(x) \psi_\ell(x+\epsilon)} = \lim_{\epsilon\rightarrow 0} \frac{1}{2\pi\i\epsilon} \varnormord{\klb*{1 + \i\frac{2\pi}{L} N_\ell \epsilon - \i\epsilon\partial_x\varphi_\ell(x)}}, %\\
%&= -\frac{1}{2\pi}\partial_x\varphi_\ell(x) + \frac{1}{L} N_\ell \numberthis \label{eqn:bos:density operator from boson identity}
\end{align*}
which is identical to \autoref{eqn:bos:density operator boson from definition}.
We can also expand to a higher order in $\epsilon$ to get
\begin{align*}
\psi_\ell^\dagger(x) \psi_\ell(x+\epsilon) &= \frac{1}{2\pi\i\epsilon} \frac{1}{1 - \frac{\pi\i\epsilon}{L} - \frac{2\pi^2\epsilon^2}{3L^2}} \normord{\powe{2\pi\i \epsilon \rho_\ell - \i \frac{1}{2!} \epsilon^2 \partial_x^2\varphi_\ell - \i \frac{1}{3!} \epsilon^3 \partial_x^3\varphi_\ell + \dots}}\\
&= \frac{1}{2\pi\i\epsilon} \klb*{1 + \frac{\pi\i\epsilon}{L} - \frac{\pi^2\epsilon^2}{3L^2}} \normord{\klb*{1 + 2\pi\i \epsilon \rho_\ell - \i \frac{1}{2!} \epsilon^2 \partial_x^2\varphi_\ell - \frac{(2\pi)^2}{2!} \epsilon^2 \rho_\ell^2}} + \mathcal{O}\kla*{\epsilon^2}
%&= \frac{1}{2\pi\i\epsilon} \normord{\klb*{1 + 2\pi\i \epsilon \rho_\ell - \i \frac{1}{2!} \epsilon^2 \partial_x^2\varphi_\ell - \frac{(2\pi)^2}{2!} \epsilon^2 \rho_\ell^2 +  + \frac{\pi\i\epsilon}{L} - \frac{2\pi^2\epsilon^2}{L} \rho_\ell - \frac{\pi^2\epsilon^2}{3L^2}}}
\end{align*}
and differentiating with respect to $\epsilon$ then yields
\begin{align*}
\psi_\ell^\dagger(x) \partial_\epsilon \psi_\ell(x+\epsilon) &= \frac{-1}{2\pi\i\epsilon^2} + \frac{1}{2\pi\i} \normord{\klb*{-\i \frac{1}{2} \partial_x^2\varphi_\ell - 2\pi^2 \rho_\ell^2 - \frac{2\pi^2}{L}\rho_\ell - \frac{\pi^2}{3L^2}}} + \mathcal{O}(\epsilon).
\end{align*}
Taking the regularized limit finally gives
\begin{align}
\varnormord{\psi_\ell^\dagger(x) \partial_x \psi_\ell(x)} = -\frac{1}{4\pi} \partial_x^2 \varphi_\ell + \pi\i \normord{\rho_\ell \kla*{\rho_\ell + \frac{1}{L}}}. \label{eqn:bos:psi_ell_partial_psi_ell_regular}
\end{align}
Note that the additional $\frac{1}{L}$ originates from using the exact commutation relation in \autoref{eqn:bos:commutator_varphi_varphi_exact} instead of its asymptotic form in \autoref{eqn:bos:commutator_varphi_varphi_tdl} and is therefore easy to miss.
This result is equivalent to the Hamiltonian derived in \cite[\ssec 7]{Bosonization_DelftSchoeller_1998} via a very different approach.\\
The commutation relations for distinct species are trivial, so the only other interesting bilinear is comparatively straightforward to bosonize, since for $\ell\neq \ell'$,
\begin{align*}
\psi_\ell^\dagger(x) \psi_{\ell'}(x') &= \normord{\powe{\i\varphi_\ell(x)}} \powe{-\i\frac{2\pi}{L} N_\ell x} \frac{\eta_\ell^\dagger \eta_{\ell'}}{L} \powe{\i \frac{2\pi}{L} N_{\ell'} x'} \normord{\powe{-\i\varphi_{\ell'}(x')}}\\
&= \frac{\eta_\ell^\dagger \eta_{\ell'}}{L} \powe{\i \frac{2\pi}{L} \klb*{N_{\ell'}x' - \kla*{N_\ell + 1}x}} \normord{\powe{\i\varphi_\ell(x) - \i\varphi_{\ell'}(x')}}. \numberthis \label{eqn:bos:psi ell dagger psi ellp neq ell}
\end{align*}
Higher order expressions are not too useful in this general form and are evaluated individually when we bosonize our model.
\myparagraph{Example: Spinless fermions}
We may now finally reap the benefits for our efforts in the previous section:
Let us demonstrate how bosonization can map fermionic interactions to quadratic terms by means of a concrete example.
For simplicity, consider the typical case of a tight-binding lattice model with hopping amplitude $t$ and lattice spacing $a$.
%Then the dispersion is $\epsilon(k) = -2t\cos(ka)$, and at half-filling the Fermi momentum is $\kf = \frac{\pi}{2a}$.
Then the dispersion is $\epsilon(k) = -2t\cos(ka)$, and for a generic Fermi momentum $\kf$, we can identify two branches by linearizing near the Fermi level (i.e. around $k=\pm \kf$).
%Linearizing near the Fermi level, i.e. around $k=\pm \kf$, we can identify two branches.
Due to the positive and negative Fermi velocity the particles are respectively referred to as right- and left-movers, and we denote them with the labels $\ell=\R,\L$ \cite{Bosonization_DelftSchoeller_1998,Giamarchi_2004}.
A fermionic field operator can then be written as
\begin{align*}
\psi(x) &= \frac{1}{\sqrt{L}} \sum_k \powe{\i kx} c_k \approx \frac{1}{\sqrt{L}} \sum_{|k-\kf|<\Lambda} \powe{\i kx} c_k + \frac{1}{\sqrt{L}}\sum_{|k+\kf|<\Lambda} \powe{\i kx} c_k
\end{align*}
where $\Lambda$ is the bandwidth.
Each of these sums corresponds to a field operator subject to \autoref{eqn:bos:bosonization identity}, and thus $\psi(x) \approx \powe{\i \kf x} \psi_\R(x) + \powe{-\i \kf x} \psi_\L(x)$ \cite{Bosonization_Miranda_2003}.
The bosonization identities read $\psi_\ell(x) = \frac{\eta_\ell}{\sqrt{L}} \normord{\powe{-\i\varphi_\ell(x)}}$ and it is customary to also introduce the fields $\varphi_\ell = \ell \phi - \theta$ \cite{Giamarchi_2004}.\\
Via \autoref{eqn:bos:psi_ell_partial_psi_ell_regular} and in the thermodynamic limit, the free Hamiltonian becomes \cite{QImpurity_review_Furusaki_2005}
\begin{align*}
H_0 &\approx \sum_\ell \sum_k c_k^\dagger \vf(\ell k - \kf) c_k = \sum_\ell \vf \int\pi \normord{\rho_\ell^2} = \frac{\vf}{2\pi} \int \normord{\klb*{\kla*{\partial_x\theta}^2 + \kla*{\partial_x\phi}^2}}.
\end{align*}%\sum_k \epsilon(k) c_k^\dagger c_k \approx \sum_\ell \sum_k c_k^\dagger (\ell k - \kf) c_k = \sum_\ell \intx{ }{ }{\varnormord{\psi_\ell(x)^\dagger (-\i\ell\partial_x) \psi(x)}}{x}
For spinless fermions, a generic interaction takes the form \cite[\ssec 2.2]{Giamarchi_2004}
\begin{align*}
H_\tint &= \int V(x-x')\rho(x)\rho(x') = \frac{1}{2L} \sum_{q,k,k'} V(q) c_{k+q}^\dagger c_{k'-q}^\dagger c_{k'} c_{k}.
\end{align*}
Within the low-energy limit, we are only interested in contributions where all particles start and end near the Fermi surface.
In the spinless case, this leads to $q\ll \kf$ and thus we can approximate the interaction as $H_\tint \approx \frac{V_0}{2L} \sum_{k,k'} n_k n_{k'}$ where $V_0 = V(q\approx 0)$.
If both particles are on the same branch, i.e. $kk'>0$, this yields terms proportional to $\rho_\ell(x) \rho_\ell(x')$.
For a sufficiently local interaction we have $x'\approx x$.
In terms of the bosonic fields this gives\footnote{For $\ell=\L$ the dispersion is decreasing, which leads to extra minus signs; see \autoref{sec:ell:bos:model_boson} for details.}
\begin{align*}
\sum_\ell \rho_\ell(x) \rho_\ell(x') &\approx \sum_\ell \frac{-\ell}{2\pi} \partial_x\varphi_\ell \frac{-\ell}{2\pi} \partial_x\varphi_\ell = \frac{1}{2\pi^2} \klb*{\kla*{\partial_x\phi}^2 + \kla*{\partial_x\theta}^2},
\end{align*}
and adding the above to $H_0$ merely renormalizes the Fermi velocity.
The other case is that the interaction couples two particles on separate branches, i.e. $kk'<0$, which results in forward scattering since the particles stay on their branch \cite[\spage 17]{Giamarchi_2004}.
We then get contributions of the form $\rho_\ell(x) \rho_{-\ell}(x')$, and combining both types of interaction results in
\begin{align*}
H_\tint &\approx V_0 \sum_{\ell,\ell'} \int \rho_\ell(x) \rho_{\ell'}(x') \approx \frac{V_0}{\pi^2} \int \kla*{\partial_x\phi(x)}^2.
\end{align*}
Consequently, the full Hamiltonian is still quadratic in the bosonic fields, and thus remains solvable.
Note that, although very basic, this example is not purely academic; variations of it are ubiquitous in research on \glspl{ll}.
% e.g. \cite{FermiLiquid_LL_review_Voit_1995,Bosonization_AndersonLocal_Giamarchi_1988,Helical_Edge_PhononBackscat_Budich_2012,Helical_RashbaImp_RG_Crepin_2012,Bosonization_SSH_Giamarchi_2023,Bosonization_DelftSchoeller_1998,Helical_SyntheticQWire_Japaridze_2013,Helical_Rashba_SpinOrbit_QSH_Stroem_2010}.
A possible application is to study disorder by amending this model with an impurity term, e.g. \cite{QImpurity_review_Furusaki_2005,QImpurity_LL_Rylands_2016,Helical_Rashba_SpinOrbit_QSH_Stroem_2010,Helical_RashbaImp_RG_Crepin_2012}.
Let us also mention that often one includes spin, which leads to additional backscattering processes, and it is then no longer possible to exactly diagonalize the Hamiltonian \cite[\spage 39]{Giamarchi_2004}.
\subsection{Correlation functions}\label{sec:bos:correlators}
In this section, we compute the correlation functions of a \gls{ll} at finite temperature.
Exact expressions for finite $L$ exist, but are rather unwieldy, see \cite{LL_FiniteT_FiniteL_Properties_Mattson_1997}.
With the thermodynamic limit in mind, we only keep contributions in the leading order of $\frac{1}{L}$.\\
Let us consider only a single species and hence drop the label $\ell$ for notational convenience.
Via \autoref{eqn:bos:psi_ell_partial_psi_ell_regular} a linear dispersion with velocity $v$ gives the Hamiltonian
\begin{align}
H &= \sum_k vk\varnormord{c_{k}^\dagger c_k} = v\intx{ }{ }{\varnormord{\psi^\dagger(x) (-\i \partial_x) \psi(x)}}{x} = \frac{v}{4\pi} \intx{ }{ }{\normord{\kla*{\partial_x\varphi(x)}^2}}{x}. \label{eqn:bos:Hamilton vk varnorm nk as 4pi norm dx varphi sqr}
\end{align}
At this point, a few remarks about the hierarchy of length scales is in order.
Our general setting is that of a lattice model at finite temperature.
The lattice spacing $a$ generates a bandwidth of order $\frac{1}{a}$ and in order to apply bosonization we add infinitely many momenta, which is valid so long as we only care about the low-energy properties of the theory.
Since the dispersion relation of any specific lattice model is generically not linear, there exists an additional bandwidth $\Lambda$ that restricts the momenta to a sufficiently small range in which a linearization is possible.
This also means that we are a priori limited to small temperatures in the sense that $v\Lambda \gg T$.
Finally, the infinitely many high-lying momenta can lead to \gls{uv} divergences in expressions that are not normal-ordered \cite[\ssec 5]{Bosonization_DelftSchoeller_1998}. %also Giamarchi (somewhere)
In practise it is often more convenient to work with such expressions though, so one usually introduces the additional artificial regularization $\alpha$ as in \autoref{eqn:bos:varphi_ell_minus_def}.\\
%Assuming that $\alpha\ll a$ this introduces an irrelevant error for the regime we care about and leads to the hierarchy $\frac{1}{\alpha}\gg \frac{1}{a} \gg \Lambda \gg \frac{T}{v}$.
%All of these conditions are somewhat qualitative, so we can already at this stage expect that obtaining quantitative results may be difficult.
As long as the slight ambiguities about the exact values of artificial parameters such as $\alpha$ or $\Lambda$ only lead to a simple change, e.g. a global prefactor, this does not present a major problem and still allows us to extract the correct physics (although quantitative predictions may be difficult). % REF
\myparagraph{Fundamental Correlator}
With $\varphi = \varphi\eplus + \varphi\eminus$ and $\varphi\eminus$ from \autoref{eqn:bos:varphi_ell_minus_def}, the Hamiltonian can be rewritten as $H\simeq\sum_{q>0} q b_q^\dagger b_q$ for $\alpha\rightarrow 0$ \cite{Bosonization_DelftSchoeller_1998}.
Since this is a simple harmonic oscillator, the imaginary time evolution with $\tau=\i t$ is governed by
\begin{align*}
\i\partial_t b_q &= -\partial_\tau b_q = \klb*{b_q,H} = vqb_q
\end{align*}
and consequently $b_q(\tau) = b_q \powe{-vq\tau}$.
Similarly $b_q^\dagger(\tau) = b_q^\dagger \powe{vq\tau}$, and writing $y=v\tau - \i x$ gives
\begin{align}
\varphi(y) &= -\sum_{q>0} \frac{\powe{-\alpha q/2}}{\sqrt{n_q}} \kla*{\powe{-qy} b_q + \powe{qy} b_q^\dagger}. \label{eqn:bos:varphi_ell_of_yalpha}
\end{align}%= -\sum_{q>0} \frac{\powe{-\alpha q/2}}{\sqrt{n_q}}  \kla*{\powe{\i qx} b_q \powe{-vq\tau} + \powe{-\i qx} b_q^\dagger \powe{vq\tau}}
Computing correlators can then be done directly on the operator level and we define
\begin{align*}
\tilde{\Phi}(y) &= \bra*{\varphi(y) \varphi(0)} = \sum_{q,q'>0} \frac{1}{\sqrt{n_qn_{q'}}} \powe{-\alpha \frac{q+q'}{2}} \bra*{\kla*{\powe{-qy} b_q + \powe{qy} b_q^\dagger} \kla*{b_{q'} + b_{q'}^\dagger}}\\
&= \sum_{q>0} \frac{1}{n_q} \powe{-\alpha q} \klb*{\powe{-qy}\bra[\big]{b_q b_q^\dagger} + \powe{qy}\bra[\big]{b_q^\dagger b_q}}= \sum_{q>0} \frac{1}{n_q} \powe{-\alpha q} \klb*{\frac{2\cosh(qy)}{\powe{\beta vq} - 1} + \powe{-qy}}.
\end{align*}
Expanding the hyperbolic cosine and the Bose function yields the same modified geometric series as in \autoref{eqn:bos:commutator_varphi_varphi_exact}, and thus
\begin{align*}
\tilde{\Phi}(y) &= \sum_{n_q\ge 1} \frac{1}{n_q} \powe{- \frac{2\pi}{L} n_q(\alpha + y)} + \sum_{s=\pm 1} \sum_{m\ge 1} \sum_{n_q\ge 1} \frac{1}{n_q} \powe{-\frac{2\pi}{L} n_q (\alpha + \beta vm + sy)}\\
&= -\log\klb*{1 - \powe{-\frac{2\pi}{L}(\alpha + y)}} - \sum_{s=\pm 1} \sum_{m\ge 1} \log\klb*{1 - \powe{-\frac{2\pi}{L} (\alpha + \beta vm + sy)}}.
\end{align*}
In the thermodynamic limit this result is plagued by an \gls{ir} divergence, so we instead consider the regularized expression
\begin{align*}
\Phi(y) &= \tilde{\Phi}(y) - \tilde{\Phi}(0) \overset{L\rightarrow\infty}{\sim} -\log\klb*{ \frac{\alpha + y}{\alpha} } - \sum_{s=\pm 1} \sum_{m\ge 1} \log\klb*{\frac{\alpha + \beta vm + sy}{\alpha + \beta vm}}\\
&= -\log\klb*{1 + \frac{y}{\alpha}} - \sum_{m\ge 1} \log \klb*{1 - \kla[\Big]{\frac{y}{\alpha + \beta vm}}^2}.
\end{align*}
Utilizing the product representation of the Gamma function one can readily verify that $\prod_{m\ge 1} \kla*{1 - \frac{x^2}{(a+m)^2}} = \frac{\Gamma(1+a)^2}{\Gamma(1+a+x)\Gamma(1+a-x)}$, and we thus obtain
\begin{align}
\Phi(y) &\overset{L\rightarrow\infty}{\sim} \log\klb*{\Gamma\kla[\Big]{1 + \frac{\alpha + y}{\beta v}} \Gamma\kla[\Big]{1 + \frac{\alpha - y}{\beta v}}} - \log\klb*{\Gamma\kla[\Big]{1 + \frac{\alpha}{\beta v}}^2 \kla*{1 + \frac{y}{\alpha}}}. \label{eqn:bos:correlator_Phi_z_alpha_exact}
\end{align}
For $\alpha\ll y$, and due to $\Gamma(1+x)\Gamma(1-x) = \pi x \csc(\pi x)$, we also find the approximation
\begin{align}
\Phi(z) &\approx \log\klb*{\frac{\pi\alpha}{\beta v} \csc\kla{z}} \label{eqn:bos:correlator_Phi_z_alpha_log_approx}
\end{align}
where $z = \frac{\pi y}{\beta v} = \frac{\pi}{\beta v} \kla*{v \tau - \i x}$ is dimensionless.
This is the form we always use, and there are four arguments for doing so.
First, one can argue that with \acrlongpl{ft} in mind the approximation should have a negligible impact for $k \ll \frac{1}{\alpha}$.
Even though we integrate over the full space, including small $x$ and $\tau$, the error made on short scales should only be visible for sufficiently large momenta.
Second, taking a more practical point of view, \autoref{eqn:bos:correlator_Phi_z_alpha_exact} makes the analytic calculation of the \acrlong{gf} in \tpspace{} intractable\footnote{And, as may be evident from the extensive \sappl, even then the calculation is rather brutal.}.
%For the sake of making progress without immediately resorting to numerical integration we are thus forced to make some approximation.
Third, and this may be the most important point, the single-particle \acrlong{gf} calculated in \autoref{sec:ell:gf:interacting_LL} is \emph{not anti-periodic}\footnote{Although it remains unclear exactly where the anti-periodicity breaks. Including the dampening $\powe{-\alpha q}$ in the Hamiltonian and evaluating the series numerically does not remedy the issue; this seems to be an even more subtle problem.} if one uses \autoref{eqn:bos:correlator_Phi_z_alpha_exact}.
%FOOTNOTE: numeric correlator calculation for exact alpha
%from mpmath import *
%beta=1.2341;v=0.551;L=1000;x=0.441;tau=0.412;alpha=1e-1
%def myfunc(n):
%    q = 2*pi/L*n; damp = exp(-alpha*q)
%    return 1/n * damp * (2*(cosh(q*v*damp*tau-1j*q*x) - 1) / (exp(beta*v*q*damp) - 1) + exp(1j*q*x - q*v*damp*tau)-1)
%resn = nsum(myfunc, [1, inf], maxterms=10**4, tol=1e-6, verbose=True, method="r")
%resa = log(pi*alpha/beta/v / sin(pi/beta/v * (v*tau - 1j*x)))
%eresn = exp(resn);eresa = exp(resa)
%print(eresn, eresa)
As discussed in \autoref{sec:gf:interaction_imaginary_time}, this anti-periodicity is a very general statement and thus, despite the mathematical inconsistency, \autoref{eqn:bos:correlator_Phi_z_alpha_log_approx} must be the physically correct result.
Finally, except for a vague mention of a result in terms of Gamma functions in \cite[\spage 22]{Bosonization_FiniteT_Bowen_2001}, the literature gives a relation akin to the approximate form, e.g. \cites[\sapp C]{Giamarchi_2004}[\ssec 6.7]{Fradkin_2013}[\ssec 3]{Bosonization_Rao_2001}[\ssec 8]{Bosonization_DelftSchoeller_1998} and \autoref{table:ell:gf:gf_literature}.
\myparagraph{Debye-Waller relation}
Later we seek to calculate correlation functions for the fermions, and due to the structure of the bosonization identity this requires expectation values of the form $\bra*{\powe{\varphi(z_1)} \dots \powe{\varphi(z_N)}}$.
One can solve this directly via path integrals \cite[\sapp C]{Giamarchi_2004}, but generalizing the problem leads to a more elegant approach.\\
Consider arbitrary linear combinations $C_j$ of bosonic operators, e.g. our $b_q$ from bosonization.
By construction their commutators are $\mathbb{C}$-valued and repeatedly applying the Baker-Campbell-Hausdorff formula gives
\begin{align*}
\prod_j \powe{C_j} &= \powe{C_1 + C_2} \powe{\frac{1}{2}\klb*{C_1, C_2}} \powe{C_3} \powe{C_4} \dots = \powe{\sum_j C_j} \powe{\frac{1}{2} \sum_{j<k} \klb*{C_j,C_k}}.
\end{align*}
The tool enabling the computation of expectation values for such complicated objects is the Debye-Waller relation
\begin{align}
\bra[\big]{\powe{C}} &= \powe{\frac{1}{2} \bra*{C^2}} \label{eqn:bos:debye_waller}
\end{align}
which holds for linear combinations $C$ of bosonic operators and, importantly, when averaging with respect to a quadratic Hamiltonian.
Two approaches for proving this result can be found in \cite[\sapp C, J]{Bosonization_DelftSchoeller_1998} and yet another in Mermin's ``one sentence proof'' \cite{DebyeWaller_1sent_Mermin_1966}.\\
By construction $\klb{C_j,C_k} = \bra{\klb{C_j,C_k}} = \bra{C_jC_k} - \bra{C_kC_j}$, and thus
\begin{align}
\bra[\Big]{\prod_j \powe{C_j}} &= \powe{\frac{1}{2} \bra[\big]{\kla[\big]{\sum_j C_j}^2}} \powe{\frac{1}{2} \sum_{j<k} \klb*{C_j, C_k}} = \powe{\frac{1}{2} \sum_j \bra{C_j^2}} \powe{\sum_{j<k} \bra{C_jC_k}}. \label{eqn:bos:expval_prod_powe_Cj}
\end{align}
For $C_j=\i A_j \varphi(z_j)$, and utilizing translational invariance, this then simplifies to
\begin{align*}
\bra[\Big]{\prod_j \powe{\i A_j \varphi(z_j)}} &= \powe{-\sum_{j<k} A_jA_k \Phi(z_j-z_k)} \powe{-\frac{1}{2} A^2 \tilde{\Phi}(0)} \numberthis\label{eqn:bos:expval_prod_powe_iAj_varphi}
\end{align*}%= \powe{-\frac{1}{2} \sum_j A_j^2\bra*{\varphi(0)^2}} \powe{-\sum_{j<k} A_jA_k \bra*{\varphi(z_j-z_k) \varphi(0)}}\\
where $A=\sum_jA_j$.
In the thermodynamic limit we have $\tilde{\Phi}(0)\rightarrow \infty$, and therefore the expectation value is only non-zero if the \emph{neutrality condition} $A=0$ is satisfied \cites[\sapp C]{Giamarchi_2004}[\sapp K]{Bosonization_DelftSchoeller_1998}.